\documentclass[10pt,a4paper]{amsart}
\usepackage[margin=19mm]{geometry}
\usepackage{amsmath,amssymb,mathtools}
\usepackage[scr=rsfs,cal=euler]{mathalfa}
\usepackage{xfrac}
\usepackage[unicode,hidelinks,bookmarks=false]{hyperref}
\newcommand{\ie}{i.e.\ }
\newcommand{\cf}{cf.\ }
\newcommand{\resp}{respectively\ }
\newcommand{\Subsection}{\subsection}
\providecommand{\nobreakdash}{\nobreak\hbox{-}}
\setlength{\emergencystretch}{2em}
\multlinegap0pt
\usepackage{amssymb}
\usepackage{enumerate}
\newtheorem{fact}{Fact}
\newtheorem{proposition}{Proposition}
\newcommand{\Ext}{\operatorname{Ext}}
\newcommand{\Min}{\operatorname{Min}}
\newcommand{\Supp}{\operatorname{Supp}}
\newcommand{\depth}{\operatorname{depth}}
\newcommand{\pd}{\operatorname{pd}}
\title{Remarks on modules of finite projective dimension}
\author{Mohsen Asgharzadeh, Elham Mahdavi}
\date{Source: arXiv:2602.09899v2 (CC BY 4.0)}
\begin{document}
\maketitle

\noindent\emph{Source and licence.} Remarks on modules of finite projective dimension. Version arXiv:2602.09899v2 (CC BY 4.0), distributed by arXiv under the Creative Commons Attribution 4.0 International licence, http://creativecommons.org/licenses/by/4.0/, as recorded in the arXiv licence metadata for this exact version and shown on the abstract page https://arxiv.org/abs/2602.09899v2. Full licence notice: \texttt{licenses/arxiv-2602.09899-CC-BY-4.0.txt}.\par\smallskip

\noindent\textit{Editorial scope.} Counted unit: the article's proposition proposition (source0.tex lines 706--711). The article's complete original proof of it is delivered with it (source0.tex lines 713--735, one \texttt{proof} environment of 23 lines). No mathematics is written, repaired or re-derived: the statement and the proof are the article's own lines, in the article's own notation, and no retained line is substituted or re-flowed.  Retained uncounted as the source-local context the proof needs: lines 649--651.  Nothing was omitted from any retained span, so the package carries no omission record.  No retained line contains a citation, so the wrapper carries no bibliography and no bibliography entry.  No retained line contains a figure environment, an image-inclusion command or drawing code, so no image is delivered and the package declares no asset dependency.  Editorial formatting only, in addition: an AMS \texttt{amsart} preamble that restates the article's own declarations and macros used by the retained lines, namely the environments \texttt{fact}, \texttt{proposition} on separate counters, together with the standard packages those lines need.  The retained environments print in this document as Fact 1, Proposition 1 in order of appearance, whereas the article prints the same items with the numbering of its own full document.  The complete native source of the article as distributed in the arXiv e-print for this exact version is bundled unchanged as \texttt{sources/194353/source0.tex}.
\begin{fact}\label{art}
	If there exists a nonzero Artinian $R$-module of finite projective dimension, then $R$ is Cohen–Macaulay.
\end{fact}

\begin{proposition}
	Let $R$ be $d$-dimensional, and let $M$ be a $1$-dimensional module with $\pd(M) < \infty$. Then 
	\(
	\dim(R/\mathfrak p) + \operatorname{ht}(\mathfrak p) = \dim(R) 
	\) for every $\mathfrak p \in \Supp(\Ext^{d-1}_R(M,R))$.
\end{proposition}

\begin{proof}
	If $\Ext^{d-1}_R(M,R)=0$, there is nothing to prove. Let $\mathfrak p \in \Supp(\Ext^{d-1}_R(M,R))$.
First, assume that $\mathfrak p = \mathfrak m$. The the equality is trivial. Otherwise $\mathfrak p \in \Min(\Supp(M))$, since $\Supp(\Ext^{d-1}(M,R)) \subseteq \Supp(M)$ and $\dim(M)=1$. Thus $\dim(M_{\mathfrak p})=0$.
	Since $\pd(M_{\mathfrak p})<\infty$, Fact \ref{art} implies $R_{\mathfrak p}$ is Cohen–Macaulay. By Auslander–Buchsbaum,
	\[
	\pd(M_{\mathfrak p}) \le \depth(R_{\mathfrak p}) = \dim(R_{\mathfrak p}) = \operatorname{ht}(\mathfrak p).
	\]
	On the other hand,
	\[
	\Ext^{d-1}_{R_{\mathfrak p}}(M_{\mathfrak p}, R_{\mathfrak p}) \neq 0
	\;\Rightarrow\;
	\pd(M_{\mathfrak p}) \ge d-1.
	\]
	Hence
	\[
	\operatorname{ht}(\mathfrak p) \ge d-1 = d - \dim(R/\mathfrak p),
	\]
	so
	\(
	\dim(R/\mathfrak p) + \operatorname{ht}(\mathfrak p) \ge d.
	\)
	The reverse inequality always holds, giving equality.
\end{proof}We now pose the following:

\end{document}
